\documentclass{econsocart}
\input{./._relpath-to-latexroot.ltx}
\usepackage{../@local/local-qe-figs-and-tables}

\begin{document}

% Add clear separation from any preceding content
\vspace{1em}
\FloatBarrier %float package: \FloatBarrier

\begin{figure}[H] %float package: [H] option
  \centering
  \caption{\pregenmark Policy effectiveness during recessions with aggregate demand effects}
  \label{fig:Policyrelrecession} % integrated doc: SST for labels
  \noindent\begin{minipage}{\textwidth}
    \centering
    \begin{subfigure}[b]{.32\linewidth}
      \centering
      % Original path: \latexroot/Code/HA-Models/FromPandemicCode/Figures/recession_Check_relrecession
      \includegraphics[width=\linewidth]{\latexroot/Code/HA-Models/FromPandemicCode/Figures/recession_Check_relrecession.pdf}
      \caption{Check IRF}
      \label{fig:recessioncheckrelrecession} % integrated doc: SST for labels
    \end{subfigure}
    \hfill%
    \begin{subfigure}[b]{.32\linewidth}
      \centering
      % Original path: \latexroot/Code/HA-Models/FromPandemicCode/Figures/recession_UI_relrecession
      \includegraphics[width=\linewidth]{\latexroot/Code/HA-Models/FromPandemicCode/Figures/recession_UI_relrecession.pdf}
      \caption{UI extn IRF}
      \label{fig:recessionuirelrecession} % integrated doc: SST for labels
    \end{subfigure}
    \hfill%
    \begin{subfigure}[b]{.32\linewidth}
      \centering
      % Original path: \latexroot/Code/HA-Models/FromPandemicCode/Figures/recession_taxcut_relrecession
      \includegraphics[width=\linewidth]{\latexroot/Code/HA-Models/FromPandemicCode/Figures/recession_taxcut_relrecession.pdf}
      \caption{tax cut IRF}
      \label{fig:recessiontaxcutrelrecession} % integrated doc: SST for labels
    \end{subfigure}
    \\[1.5em]
    \begin{subfigure}[b]{.32\linewidth}
      \centering
      % Original path: \latexroot/Code/HA-Models/FromPandemicCode/Figures/Cumulative_multiplier_Check
      \includegraphics[width=\linewidth]{\latexroot/Code/HA-Models/FromPandemicCode/Figures/Cumulative_multiplier_Check.pdf}
      \caption{Check multiplier}
      \label{fig:recessioncheckrelrecession_Mult} % integrated doc: SST for labels
    \end{subfigure}
    \hfill%
    \begin{subfigure}[b]{.32\linewidth}
      \centering
      % Original path: \latexroot/Code/HA-Models/FromPandemicCode/Figures/Cumulative_multiplier_UI
      \includegraphics[width=\linewidth]{\latexroot/Code/HA-Models/FromPandemicCode/Figures/Cumulative_multiplier_UI.pdf}
      \caption{UI extn multiplier}
      \label{fig:recessionuirelrecession_Mult} % integrated doc: SST for labels
    \end{subfigure}
    \hfill%
    \begin{subfigure}[b]{.32\linewidth}
      \centering
      % Original path: \latexroot/Code/HA-Models/FromPandemicCode/Figures/Cumulative_multiplier_TaxCut
      \includegraphics[width=\linewidth]{\latexroot/Code/HA-Models/FromPandemicCode/Figures/Cumulative_multiplier_TaxCut.pdf}
      \caption{tax cut multiplier}
      \label{fig:recessiontaxcutrelrecession_Mult} % integrated doc: SST for labels
    \end{subfigure}
  \end{minipage}
\end{figure}
\noindent\parbox{\textwidth}{\footnotesize
  \textbf{Note}: This figure compares policy effectiveness during recessions (Section~\ref{sec:recessions}).
  Recession periods are characterized by higher unemployment rates and increased economic uncertainty.
  The model demonstrates that UI extensions become particularly effective during recessions due to
  better targeting of high-MPC households. Stimulus checks maintain effectiveness but with diminished
  relative performance compared to UI extensions. Payroll tax cuts show the least effectiveness
  across all economic conditions, confirming the robustness of the main policy rankings.
}

% Add clear separation from following content
\vspace{1em}
\FloatBarrier %float package: \FloatBarrier


\end{document}
